\import{algebra/ring/congruence.tex}
\import{algebra/ring/homomorphism.tex}
\import{algebra/group/quotient.tex}
\import{relation/quotient-lift.tex}

\section{Quotient rings}

Let $I$ be a two-sided ideal of a ring $R$.  The additive structure of
$R/I$ is the existing quotient of the additive group.  Multiplication is
descended directly through the same additive-coset relation, without choosing
representatives, and the quotient unit is literally the projection of
$1_R$.

\subsection{The descended multiplication graph}

\begin{definition}\label{quotient_ring_representative_multiplication}
    $\quotientRingRepresentativeMultiplication{R}{I}=
        \groupQuotientProjection{\AdditiveGroup{R}}{I}
            \circ\ringMul[R]$.
\end{definition}

\begin{proposition}\label{quotient_ring_representative_multiplication_type}
    Suppose $I$ is a two-sided ideal of $R$.
    Let $G=\AdditiveGroup{R}$.
    Then $\quotientRingRepresentativeMultiplication{R}{I}\in
        \funs{\carrier[R]\times\carrier[R]}{
            \groupQuotientCarrier{G}{I}}$.
\end{proposition}
\begin{proof}
    $G$ is a commutative group by
        \obv{additive_group_is_exact_commutative_group}.
    $\subgroup{I}{G}$ by
        \obv{two_sided_ideal,left_ideal_additive_subgroup}.
    Hence $\normalSubgroup{I}{G}$ by
        \obv{commutative_group_subgroup_is_normal}.
    $\ringMul[R]\in
        \funs{\carrier[R]\times\carrier[R]}{\carrier[R]}$ by
        \obv{ring_mul_graph}.
    $\groupQuotientProjection{G}{I}\in
        \funs{\carrier[G]}{\groupQuotientCarrier{G}{I}}$ by
        \obv{group_quotient_projection_type}.
    $\carrier[G]=\carrier[R]$ by \auto{additive_group_fields}.
    Follows by
        \auto{quotient_ring_representative_multiplication,funs_circ}.
\end{proof}

\begin{proposition}\label{quotient_ring_representative_multiplication_apply}
    Suppose $I$ is a two-sided ideal of $R$.
    Let $a,b\in\carrier[R]$.
    Let $G=\AdditiveGroup{R}$ and $E=\subgroupCosetRelation{G}{I}$.
    Then
        $\apply{\quotientRingRepresentativeMultiplication{R}{I}}{(a,b)}
            =\equivalenceClass{E}{\ringMul[R](a,b)}$.
\end{proposition}
\begin{proof}
    $R$ is a ring by \obv{two_sided_ideal_ring}.
    $G$ is a commutative group by
        \obv{additive_group_is_exact_commutative_group}.
    $\subgroup{I}{G}$ by
        \obv{two_sided_ideal,left_ideal_additive_subgroup}.
    Hence $\normalSubgroup{I}{G}$ by
        \obv{commutative_group_subgroup_is_normal}.
    Let $q=\groupQuotientProjection{G}{I}$.
    $\ringMul[R]$ is a function and
        $\dom{\ringMul[R]}=\carrier[R]\times\carrier[R]$ by
        \auto{ring_mul_graph,funs_elim}.
    $q$ is a function and $\dom{q}=\carrier[G]$ by
        \auto{group_quotient_projection_type,funs_elim}.
    $\carrier[G]=\carrier[R]$ by \auto{additive_group_fields}.
    $\ran{\ringMul[R]}\subseteq\carrier[R]$ by
        \auto{ring_mul_graph,funs_ran}.
    Hence $\ran{\ringMul[R]}\subseteq\dom{q}$.
    Hence $q$ is composable with $\ringMul[R]$.
    $\ringMul[R](a,b)\in\carrier[R]$ by \obv{ring_mul_type}.
    $(q\circ\ringMul[R])((a,b))=q(\ringMul[R](a,b))$ by
        \auto{circ_apply,times_tuple_intro}.
    $q(\ringMul[R](a,b))
        =\equivalenceClass{E}{\ringMul[R](a,b)}$ by
        \auto{group_quotient_projection_apply}.
    Follows by \auto{quotient_ring_representative_multiplication}.
\end{proof}

\begin{proposition}\label{quotient_ring_representative_multiplication_compatible}
    Suppose $I$ is a two-sided ideal of $R$.
    Let $G=\AdditiveGroup{R}$ and $E=\subgroupCosetRelation{G}{I}$.
    Then
        $\quotientBinaryCompatible
            {\quotientRingRepresentativeMultiplication{R}{I}}
            {\carrier[R]}{E}{\carrier[R]}{E}$.
\end{proposition}
\begin{proof}
    $\subgroup{I}{G}$ by
        \obv{two_sided_ideal,left_ideal_additive_subgroup}.
    $E$ is an equivalence on $\carrier[G]$ by
        \auto{subgroup_coset_relation_equivalence}.
    $\carrier[G]=\carrier[R]$ by \auto{additive_group_fields}.
    Show that for all $a,a'\in\carrier[R]$ we have
        for all $b,b'\in\carrier[R]$ such that
            $a\mathrel{E}a'$ and $b\mathrel{E}b'$ we have
            $\apply{\quotientRingRepresentativeMultiplication{R}{I}}{
                (a,b)}
            =\apply{\quotientRingRepresentativeMultiplication{R}{I}}{
                (a',b')}$.
    \begin{subproof}
        Fix $a,a'\in\carrier[R]$.
        Show that for all $b,b'\in\carrier[R]$ such that
            $a\mathrel{E}a'$ and $b\mathrel{E}b'$ we have
            $\apply{\quotientRingRepresentativeMultiplication{R}{I}}{
                (a,b)}
            =\apply{\quotientRingRepresentativeMultiplication{R}{I}}{
                (a',b')}$.
        \begin{subproof}
            Fix $b,b'\in\carrier[R]$.
            Assume $a\mathrel{E}a'$ and $b\mathrel{E}b'$.
            $\ringMul[R](a,b)\mathrel{E}\ringMul[R](a',b')$ by
                \obv{ideal_coset_relation_multiplication_congruent}.
            $\ringMul[R](a,b),\ringMul[R](a',b')\in\carrier[R]$ by
                \obv{ring_mul_type}.
            Hence
                $\equivalenceClass{E}{\ringMul[R](a,b)}
                    =\equivalenceClass{E}{\ringMul[R](a',b')}$ by
                \obv{equiv_implies_equivanlence_classes_eq}.
            Follows by
                \obv{quotient_ring_representative_multiplication_apply}.
        \end{subproof}
    \end{subproof}
    Follows by \obv{quotient_binary_compatible}.
\end{proof}

\begin{definition}\label{quotient_ring_multiplication_graph}
    $\quotientRingMultiplicationGraph{R}{I}=
        \quotientBinaryDescent
            {\quotientRingRepresentativeMultiplication{R}{I}}
            {\carrier[R]}
            {\subgroupCosetRelation{\AdditiveGroup{R}}{I}}
            {\carrier[R]}
            {\subgroupCosetRelation{\AdditiveGroup{R}}{I}}$.
\end{definition}

\begin{proposition}\label{quotient_ring_multiplication_graph_type}
    Suppose $I$ is a two-sided ideal of $R$.
    Let $G=\AdditiveGroup{R}$.
    Then $\quotientRingMultiplicationGraph{R}{I}\in
        \funs{
            \groupQuotientCarrier{G}{I}\times
            \groupQuotientCarrier{G}{I}}
            {\groupQuotientCarrier{G}{I}}$.
\end{proposition}
\begin{proof}
    Let $E=\subgroupCosetRelation{G}{I}$ and
        $f=\quotientRingRepresentativeMultiplication{R}{I}$.
    $\subgroup{I}{G}$ by
        \obv{two_sided_ideal,left_ideal_additive_subgroup}.
    $E$ is an equivalence on $\carrier[G]$ by
        \auto{subgroup_coset_relation_equivalence}.
    $\carrier[G]=\carrier[R]$ by \auto{additive_group_fields}.
    $f\in\funs{\carrier[R]\times\carrier[R]}{
        \groupQuotientCarrier{G}{I}}$ by
        \auto{quotient_ring_representative_multiplication_type}.
    $\quotientBinaryCompatible{f}{\carrier[R]}{E}{\carrier[R]}{E}$ by
        \auto{quotient_ring_representative_multiplication_compatible}.
    Follows by
        \auto{quotient_ring_multiplication_graph,
            group_quotient_carrier,quotient_binary_descent_type}.
\end{proof}

\begin{proposition}\label{quotient_ring_multiplication_graph_apply}
    Suppose $I$ is a two-sided ideal of $R$.
    Let $a,b\in\carrier[R]$.
    Let $G=\AdditiveGroup{R}$ and $E=\subgroupCosetRelation{G}{I}$.
    Then
        $\apply{\quotientRingMultiplicationGraph{R}{I}}{
            (\equivalenceClass{E}{a},\equivalenceClass{E}{b})}
        =\equivalenceClass{E}{\ringMul[R](a,b)}$.
\end{proposition}
\begin{proof}
    Let $f=\quotientRingRepresentativeMultiplication{R}{I}$.
    $\subgroup{I}{G}$ by
        \obv{two_sided_ideal,left_ideal_additive_subgroup}.
    $E$ is an equivalence on $\carrier[G]$ by
        \auto{subgroup_coset_relation_equivalence}.
    $\carrier[G]=\carrier[R]$ by \auto{additive_group_fields}.
    $f\in\funs{\carrier[R]\times\carrier[R]}{
        \groupQuotientCarrier{G}{I}}$ by
        \auto{quotient_ring_representative_multiplication_type}.
    $\quotientBinaryCompatible{f}{\carrier[R]}{E}{\carrier[R]}{E}$ by
        \auto{quotient_ring_representative_multiplication_compatible}.
    $\apply{\quotientBinaryDescent{f}{\carrier[R]}{E}{\carrier[R]}{E}}{
        (\equivalenceClass{E}{a},\equivalenceClass{E}{b})}
        =f((a,b))$ by \obv{quotient_binary_descent_apply}.
    $f((a,b))=\equivalenceClass{E}{\ringMul[R](a,b)}$ by
        \auto{quotient_ring_representative_multiplication_apply}.
    Follows by \auto{quotient_ring_multiplication_graph}.
\end{proof}

\subsection{The exact quotient ring}

\begin{definition}\label{quotient_ring}
    $\QuotientRing{R}{I}=\structure[ring]{
        \carrier=\carrier[\quotientGroup{\AdditiveGroup{R}}{I}],
        \ringAdd=\mul[\quotientGroup{\AdditiveGroup{R}}{I}],
        \ringZero=\neutral[\quotientGroup{\AdditiveGroup{R}}{I}],
        \ringNeg=\groupInv[\quotientGroup{\AdditiveGroup{R}}{I}],
        \ringMul=\quotientRingMultiplicationGraph{R}{I},
        \ringOne=
            \apply{\groupQuotientProjection{\AdditiveGroup{R}}{I}}{
                \ringOne[R]}
    }$.
\end{definition}

\begin{proposition}\label{quotient_ring_fields}
    Let $G=\AdditiveGroup{R}$ and $Q=\quotientGroup{G}{I}$.
    Let $S=\QuotientRing{R}{I}$.
    Then $\carrier[S]=\carrier[Q]$ and
        $\ringAdd[S]=\mul[Q]$ and
        $\ringZero[S]=\neutral[Q]$ and
        $\ringNeg[S]=\groupInv[Q]$ and
        $\ringMul[S]=\quotientRingMultiplicationGraph{R}{I}$ and
        $\ringOne[S]
            =\apply{\groupQuotientProjection{G}{I}}{\ringOne[R]}$.
\end{proposition}
\begin{proof}
    Follows by
        \auto{quotient_ring,
            ringaggregateprojection_carrier,ringaggregateprojection_ringAdd,
            ringaggregateprojection_ringZero,ringaggregateprojection_ringNeg,
            ringaggregateprojection_ringMul,ringaggregateprojection_ringOne}.
\end{proof}

\begin{proposition}\label{quotient_ring_carrier}
    $\carrier[\QuotientRing{R}{I}]
        =\groupQuotientCarrier{\AdditiveGroup{R}}{I}$.
\end{proposition}
\begin{proof}
    Follows by \auto{quotient_ring_fields,quotient_group_carrier}.
\end{proof}

\begin{proposition}\label{quotient_ring_zero}
    Let $G=\AdditiveGroup{R}$ and $E=\subgroupCosetRelation{G}{I}$.
    Then $\ringZero[\QuotientRing{R}{I}]
        =\equivalenceClass{E}{\ringZero[R]}$.
\end{proposition}
\begin{proof}
    Follows by
        \auto{quotient_ring_fields,quotient_group_neutral,
            additive_group_fields}.
\end{proof}

\begin{proposition}\label{quotient_ring_one}
    Suppose $I$ is a two-sided ideal of $R$.
    Let $G=\AdditiveGroup{R}$ and $E=\subgroupCosetRelation{G}{I}$.
    Then $\ringOne[\QuotientRing{R}{I}]
        =\equivalenceClass{E}{\ringOne[R]}$.
\end{proposition}
\begin{proof}
    $R$ is a ring by \obv{two_sided_ideal_ring}.
    $G$ is a commutative group by
        \obv{additive_group_is_exact_commutative_group}.
    $\subgroup{I}{G}$ by
        \obv{two_sided_ideal,left_ideal_additive_subgroup}.
    Hence $\normalSubgroup{I}{G}$ by
        \obv{commutative_group_subgroup_is_normal}.
    $\ringOne[R]\in\carrier[G]$ by
        \auto{ring_one_type,additive_group_fields}.
    Follows by
        \auto{quotient_ring_fields,group_quotient_projection_apply}.
\end{proof}

\begin{proposition}\label{quotient_ring_add_apply}
    Suppose $I$ is a two-sided ideal of $R$.
    Let $a,b\in\carrier[R]$.
    Let $G=\AdditiveGroup{R}$ and $E=\subgroupCosetRelation{G}{I}$.
    Then
        $\ringAdd[\QuotientRing{R}{I}]
            (\equivalenceClass{E}{a},\equivalenceClass{E}{b})
        =\equivalenceClass{E}{\ringAdd[R](a,b)}$.
\end{proposition}
\begin{proof}
    $G$ is a commutative group by
        \obv{two_sided_ideal_ring,
            additive_group_is_exact_commutative_group}.
    $\subgroup{I}{G}$ by
        \obv{two_sided_ideal,left_ideal_additive_subgroup}.
    Hence $\normalSubgroup{I}{G}$ by
        \obv{commutative_group_subgroup_is_normal}.
    $a,b\in\carrier[G]$ by \auto{additive_group_fields}.
    $\ringAdd[R](a,b)=\mul[G](a,b)$ by
        \auto{additive_group_fields}.
    Follows by
        \auto{quotient_ring_fields,quotient_group_mul_apply}.
\end{proof}

\begin{proposition}\label{quotient_ring_neg_apply}
    Suppose $I$ is a two-sided ideal of $R$.
    Let $a\in\carrier[R]$.
    Let $G=\AdditiveGroup{R}$ and $E=\subgroupCosetRelation{G}{I}$.
    Then
        $\ringNeg[\QuotientRing{R}{I}](\equivalenceClass{E}{a})
        =\equivalenceClass{E}{\ringNeg[R](a)}$.
\end{proposition}
\begin{proof}
    $G$ is a commutative group by
        \obv{two_sided_ideal_ring,
            additive_group_is_exact_commutative_group}.
    $\subgroup{I}{G}$ by
        \obv{two_sided_ideal,left_ideal_additive_subgroup}.
    Hence $\normalSubgroup{I}{G}$ by
        \obv{commutative_group_subgroup_is_normal}.
    $a\in\carrier[G]$ by \auto{additive_group_fields}.
    $\ringNeg[R](a)=\groupInv[G](a)$ by
        \auto{additive_group_fields}.
    Follows by
        \auto{quotient_ring_fields,quotient_group_inv_apply}.
\end{proof}

\begin{proposition}\label{quotient_ring_mul_apply}
    Suppose $I$ is a two-sided ideal of $R$.
    Let $a,b\in\carrier[R]$.
    Let $G=\AdditiveGroup{R}$ and $E=\subgroupCosetRelation{G}{I}$.
    Then
        $\ringMul[\QuotientRing{R}{I}]
            (\equivalenceClass{E}{a},\equivalenceClass{E}{b})
        =\equivalenceClass{E}{\ringMul[R](a,b)}$.
\end{proposition}
\begin{proof}
    Follows by
        \auto{quotient_ring_fields,
            quotient_ring_multiplication_graph_apply}.
\end{proof}

\begin{theorem}\label{quotient_ring_is_exact}
    Suppose $I$ is a two-sided ideal of $R$.
    Then $\QuotientRing{R}{I}$ is an exact ring.
\end{theorem}
\begin{proof}
    Let $G=\AdditiveGroup{R}$ and $E=\subgroupCosetRelation{G}{I}$.
    Let $Q=\quotientGroup{G}{I}$ and $S=\QuotientRing{R}{I}$.
    $R$ is a ring by \obv{two_sided_ideal_ring}.
    $G$ is a commutative group by
        \obv{additive_group_is_exact_commutative_group}.
    Hence $G$ is a group by \obv{commutative_group_is_group}.
    $\subgroup{I}{G}$ by
        \obv{two_sided_ideal,left_ideal_additive_subgroup}.
    Hence $\normalSubgroup{I}{G}$ by
        \obv{commutative_group_subgroup_is_normal}.
    $Q$ is an exact commutative group by
        \auto{commutative_quotient_group_is_exact}.
    Hence $Q$ is a group and $Q$ is a monoid and
        $Q$ is a unital magma by
        \obv{commutative_group_is_group,group_is_monoid,
            monoid_is_unital_magma}.
    $\carrier[S]=\carrier[Q]$ and
        $\ringAdd[S]=\mul[Q]$ and
        $\ringZero[S]=\neutral[Q]$ and
        $\ringNeg[S]=\groupInv[Q]$ and
        $\ringMul[S]=\quotientRingMultiplicationGraph{R}{I}$ and
        $\ringOne[S]
            =\apply{\groupQuotientProjection{G}{I}}{\ringOne[R]}$ by
        \obv{quotient_ring_fields}.
    $\ringAdd[S]\in
        \funs{\carrier[S]\times\carrier[S]}{\carrier[S]}$ by
        \auto{monoid_is_magma,magma_mul_graph}.
    $\ringZero[S]\in\carrier[S]$ by \auto{unitalmagma_type}.
    $\ringNeg[S]\in\funs{\carrier[S]}{\carrier[S]}$ by
        \auto{group_inverse_graph}.
    $\ringMul[S]\in
        \funs{\carrier[S]\times\carrier[S]}{\carrier[S]}$ by
        \auto{quotient_ring_multiplication_graph_type,
            quotient_group_carrier}.
    $\groupQuotientProjection{G}{I}\in
        \funs{\carrier[G]}{\carrier[Q]}$ by
        \auto{group_quotient_projection_type,quotient_group_carrier}.
    $\ringOne[R]\in\carrier[G]$ by
        \auto{ring_one_type,additive_group_fields}.
    Hence $\ringOne[S]\in\carrier[S]$ by \auto{funs_type_apply}.
    For all $U,V,W\in\carrier[S]$ we have
        $\ringAdd[S](U,\ringAdd[S](V,W))
            =\ringAdd[S](\ringAdd[S](U,V),W)$ by
        \auto{monoid_assoc}.
    For all $U,V\in\carrier[S]$ we have
        $\ringAdd[S](U,V)=\ringAdd[S](V,U)$ by
        \auto{commutative_group_commutative}.
    For every $U\in\carrier[S]$ we have
        $\ringAdd[S](\ringZero[S],U)=U$ by
        \auto{unitalmagma_left}.
    For every $U\in\carrier[S]$ we have
        $\ringAdd[S](\ringNeg[S](U),U)=\ringZero[S]$ by
        \auto{group_inverse_left}.
    Show that for all $U,V,W\in\carrier[S]$ we have
        $\ringMul[S](U,\ringMul[S](V,W))
            =\ringMul[S](\ringMul[S](U,V),W)$.
    \begin{subproof}
        Fix $U,V,W\in\carrier[S]$.
        Take $a\in\carrier[G]$ such that
            $U=\equivalenceClass{E}{a}$ by
            \auto{quotient_ring_carrier,group_quotient_carrier,quotient}.
        Take $b\in\carrier[G]$ such that
            $V=\equivalenceClass{E}{b}$ by
            \auto{quotient_ring_carrier,group_quotient_carrier,quotient}.
        Take $c\in\carrier[G]$ such that
            $W=\equivalenceClass{E}{c}$ by
            \auto{quotient_ring_carrier,group_quotient_carrier,quotient}.
        Hence $a,b,c\in\carrier[R]$ by \auto{additive_group_fields}.
        $\ringMul[R](a,b),\ringMul[R](b,c)\in\carrier[R]$ by
            \obv{ring_mul_type}.
        $\ringMul[S](V,W)=\equivalenceClass{E}{\ringMul[R](b,c)}$ and
            $\ringMul[S](U,V)=\equivalenceClass{E}{\ringMul[R](a,b)}$ by
            \auto{quotient_ring_mul_apply}.
        Hence
            $\ringMul[S](U,\ringMul[S](V,W))
                =\equivalenceClass{E}{
                    \ringMul[R](a,\ringMul[R](b,c))}$ and
            $\ringMul[S](\ringMul[S](U,V),W)
                =\equivalenceClass{E}{
                    \ringMul[R](\ringMul[R](a,b),c)}$ by
            \auto{quotient_ring_mul_apply}.
        Follows by \obv{ring_mul_associative}.
    \end{subproof}
    Show that for every $U\in\carrier[S]$ we have
        $\ringMul[S](\ringOne[S],U)=U$ and
        $\ringMul[S](U,\ringOne[S])=U$.
    \begin{subproof}
        Fix $U\in\carrier[S]$.
        Take $a\in\carrier[G]$ such that
            $U=\equivalenceClass{E}{a}$ by
            \auto{quotient_ring_carrier,group_quotient_carrier,quotient}.
        Hence $a\in\carrier[R]$ by \auto{additive_group_fields}.
        $\ringOne[R]\in\carrier[R]$ by \obv{ring_one_type}.
        $\ringOne[S]=\equivalenceClass{E}{\ringOne[R]}$ by
            \auto{quotient_ring_one}.
        $\ringMul[S](\ringOne[S],U)
            =\equivalenceClass{E}{\ringMul[R](\ringOne[R],a)}$ by
            \auto{quotient_ring_mul_apply}.
        $\ringMul[S](U,\ringOne[S])
            =\equivalenceClass{E}{\ringMul[R](a,\ringOne[R])}$ by
            \auto{quotient_ring_mul_apply}.
        Follows by \auto{ring_mul_one_left,ring_mul_one_right}.
    \end{subproof}
    Show that for all $U,V,W\in\carrier[S]$ we have
        $\ringMul[S](U,\ringAdd[S](V,W))
            =\ringAdd[S](\ringMul[S](U,V),\ringMul[S](U,W))$.
    \begin{subproof}
        Fix $U,V,W\in\carrier[S]$.
        Take $a\in\carrier[G]$ such that
            $U=\equivalenceClass{E}{a}$ by
            \auto{quotient_ring_carrier,group_quotient_carrier,quotient}.
        Take $b\in\carrier[G]$ such that
            $V=\equivalenceClass{E}{b}$ by
            \auto{quotient_ring_carrier,group_quotient_carrier,quotient}.
        Take $c\in\carrier[G]$ such that
            $W=\equivalenceClass{E}{c}$ by
            \auto{quotient_ring_carrier,group_quotient_carrier,quotient}.
        Hence $a,b,c\in\carrier[R]$ by \auto{additive_group_fields}.
        $\ringAdd[R](b,c),\ringMul[R](a,b),\ringMul[R](a,c)
            \in\carrier[R]$ by \obv{ring_add_type,ring_mul_type}.
        $\ringAdd[S](V,W)=\equivalenceClass{E}{\ringAdd[R](b,c)}$ by
            \auto{quotient_ring_add_apply}.
        $\ringMul[S](U,V)=\equivalenceClass{E}{\ringMul[R](a,b)}$ and
            $\ringMul[S](U,W)=\equivalenceClass{E}{\ringMul[R](a,c)}$ by
            \auto{quotient_ring_mul_apply}.
        Hence
            $\ringMul[S](U,\ringAdd[S](V,W))
                =\equivalenceClass{E}{
                    \ringMul[R](a,\ringAdd[R](b,c))}$ by
            \auto{quotient_ring_mul_apply}.
        Hence
            $\ringAdd[S](\ringMul[S](U,V),\ringMul[S](U,W))
                =\equivalenceClass{E}{
                    \ringAdd[R](\ringMul[R](a,b),\ringMul[R](a,c))}$ by
            \auto{quotient_ring_add_apply}.
        Follows by \auto{ring_distributive_left}.
    \end{subproof}
    Show that for all $U,V,W\in\carrier[S]$ we have
        $\ringMul[S](\ringAdd[S](U,V),W)
            =\ringAdd[S](\ringMul[S](U,W),\ringMul[S](V,W))$.
    \begin{subproof}
        Fix $U,V,W\in\carrier[S]$.
        Take $a\in\carrier[G]$ such that
            $U=\equivalenceClass{E}{a}$ by
            \auto{quotient_ring_carrier,group_quotient_carrier,quotient}.
        Take $b\in\carrier[G]$ such that
            $V=\equivalenceClass{E}{b}$ by
            \auto{quotient_ring_carrier,group_quotient_carrier,quotient}.
        Take $c\in\carrier[G]$ such that
            $W=\equivalenceClass{E}{c}$ by
            \auto{quotient_ring_carrier,group_quotient_carrier,quotient}.
        Hence $a,b,c\in\carrier[R]$ by \auto{additive_group_fields}.
        $\ringAdd[R](a,b),\ringMul[R](a,c),\ringMul[R](b,c)
            \in\carrier[R]$ by \obv{ring_add_type,ring_mul_type}.
        $\ringAdd[S](U,V)=\equivalenceClass{E}{\ringAdd[R](a,b)}$ by
            \auto{quotient_ring_add_apply}.
        $\ringMul[S](U,W)=\equivalenceClass{E}{\ringMul[R](a,c)}$ and
            $\ringMul[S](V,W)=\equivalenceClass{E}{\ringMul[R](b,c)}$ by
            \auto{quotient_ring_mul_apply}.
        Hence
            $\ringMul[S](\ringAdd[S](U,V),W)
                =\equivalenceClass{E}{
                    \ringMul[R](\ringAdd[R](a,b),c)}$ by
            \auto{quotient_ring_mul_apply}.
        Hence
            $\ringAdd[S](\ringMul[S](U,W),\ringMul[S](V,W))
                =\equivalenceClass{E}{
                    \ringAdd[R](\ringMul[R](a,c),\ringMul[R](b,c))}$ by
            \auto{quotient_ring_add_apply}.
        Follows by \auto{ring_distributive_right}.
    \end{subproof}
    Hence $S$ is a ring by
        \auto{quotient_ring,ringaggregateexact,ring}.
    Follows by \auto{quotient_ring,ringaggregateexact}.
\end{proof}

\subsection{The canonical projection}

\begin{theorem}\label{quotient_ring_additive_group}
    Suppose $I$ is a two-sided ideal of $R$.
    Then
        $\AdditiveGroup{\QuotientRing{R}{I}}
            =\quotientGroup{\AdditiveGroup{R}}{I}$.
\end{theorem}
\begin{proof}
    Let $G=\AdditiveGroup{R}$ and $Q=\quotientGroup{G}{I}$.
    Let $S=\QuotientRing{R}{I}$ and $H=\AdditiveGroup{S}$.
    $Q$ is an exact commutative group by
        \auto{two_sided_ideal,left_ideal_additive_subgroup,
            additive_group_is_exact_commutative_group,
            commutative_group_subgroup_is_normal,
            commutative_quotient_group_is_exact}.
    $S$ is an exact ring by \auto{quotient_ring_is_exact}.
    Hence $H$ is an exact commutative group by
        \auto{additive_group_is_exact_commutative_group}.
    $\carrier[H]=\carrier[S]=\carrier[Q]$ by
        \auto{additive_group_fields,quotient_ring_fields}.
    Hence $\carrier[H]=\carrier[Q]$.
    $\mul[H]=\ringAdd[S]=\mul[Q]$ by
        \auto{additive_group_fields,quotient_ring_fields}.
    Hence $\mul[H]=\mul[Q]$.
    $\neutral[H]=\ringZero[S]=\neutral[Q]$ by
        \auto{additive_group_fields,quotient_ring_fields}.
    Hence $\neutral[H]=\neutral[Q]$.
    $\groupInv[H]=\ringNeg[S]=\groupInv[Q]$ by
        \auto{additive_group_fields,quotient_ring_fields}.
    Hence $\groupInv[H]=\groupInv[Q]$.
    Hence $H=Q$ by structure extensionality.
\end{proof}

\begin{theorem}\label{quotient_ring_projection_hom}
    Suppose $I$ is a two-sided ideal of $R$.
    Then $\groupQuotientProjection{\AdditiveGroup{R}}{I}$ is a ring
        homomorphism from $R$ to $\QuotientRing{R}{I}$.
\end{theorem}
\begin{proof}
    Let $G=\AdditiveGroup{R}$ and $Q=\quotientGroup{G}{I}$.
    Let $S=\QuotientRing{R}{I}$ and $q=\groupQuotientProjection{G}{I}$.
    Let $M=\MultiplicativeMonoid{R}$ and
        $N=\MultiplicativeMonoid{S}$.
    $R$ is a ring by \obv{two_sided_ideal_ring}.
    $G$ is a commutative group by
        \obv{additive_group_is_exact_commutative_group}.
    $\subgroup{I}{G}$ by
        \auto{two_sided_ideal,left_ideal_additive_subgroup}.
    Hence $\normalSubgroup{I}{G}$ by
        \obv{commutative_group_subgroup_is_normal}.
    $S$ is an exact ring by \auto{quotient_ring_is_exact}.
    $\AdditiveGroup{S}=Q$ by \auto{quotient_ring_additive_group}.
    $\groupHom{q}{G}{Q}$ by \auto{group_quotient_projection_hom}.
    Hence $\groupHom{q}{G}{\AdditiveGroup{S}}$ by \auto.
    $M$ is a monoid and $N$ is a monoid by
        \auto{multiplicative_monoid_is_exact_monoid}.
    $q\in\funs{\carrier[M]}{\carrier[N]}$ by
        \auto{group_quotient_projection_type,quotient_ring_carrier,
            additive_group_fields,multiplicative_monoid_fields}.
    Show that for all $a,b\in\carrier[M]$ we have
        $q(\mul[M](a,b))=\mul[N](q(a),q(b))$.
    \begin{subproof}
        Fix $a,b\in\carrier[M]$.
        Hence $a,b\in\carrier[R]$ by
            \auto{multiplicative_monoid_fields}.
        $\mul[M](a,b)=\ringMul[R](a,b)$ and
            $\mul[N](q(a),q(b))=\ringMul[S](q(a),q(b))$ by
            \auto{multiplicative_monoid_fields}.
        $q(a)=\equivalenceClass{\subgroupCosetRelation{G}{I}}{a}$ and
            $q(b)=\equivalenceClass{\subgroupCosetRelation{G}{I}}{b}$ and
            $q(\ringMul[R](a,b))=
                \equivalenceClass{\subgroupCosetRelation{G}{I}}{
                    \ringMul[R](a,b)}$ by
            \auto{ring_mul_type,additive_group_fields,
                group_quotient_projection_apply}.
        $\ringMul[S](q(a),q(b))=
            \equivalenceClass{\subgroupCosetRelation{G}{I}}{
                \ringMul[R](a,b)}$ by
            \auto{quotient_ring_mul_apply}.
    \end{subproof}
    Hence $q$ is a magma homomorphism from $M$ to $N$ by
        \obv{monoid_is_magma,magma_hom}.
    $q(\neutral[M])=\neutral[N]$ by
        \auto{multiplicative_monoid_fields,quotient_ring_fields}.
    Hence $q$ is a monoid homomorphism from $M$ to $N$ by
        \obv{monoid_hom}.
    Follows by \auto{ring_hom}.
\end{proof}

\begin{corollary}\label{quotient_ring_projection_eq_iff_ideal_congruence}
    Suppose $I$ is a two-sided ideal of $R$.
    Let $x,y\in\carrier[R]$.
    Let $q=\groupQuotientProjection{\AdditiveGroup{R}}{I}$.
    Then $q(x)=q(y)$ iff $x\mathrel{\idealCongruence{R}{I}}y$.
\end{corollary}
\begin{proof}
    Follows by
        \auto{ideal_congruence_quotient_projection_eq_iff}.
\end{proof}

\begin{theorem}\label{quotient_ring_projection_surjective}
    Suppose $I$ is a two-sided ideal of $R$.
    Then $\groupQuotientProjection{\AdditiveGroup{R}}{I}$ is a ring
        surjection from $R$ to $\QuotientRing{R}{I}$.
\end{theorem}
\begin{proof}
    Let $G=\AdditiveGroup{R}$ and $q=\groupQuotientProjection{G}{I}$.
    $G$ is a commutative group by
        \obv{two_sided_ideal_ring,
            additive_group_is_exact_commutative_group}.
    $\subgroup{I}{G}$ by
        \auto{two_sided_ideal,left_ideal_additive_subgroup}.
    Hence $\normalSubgroup{I}{G}$ by
        \obv{commutative_group_subgroup_is_normal}.
    $q\in\surj{\carrier[G]}{\groupQuotientCarrier{G}{I}}$ by
        \auto{group_quotient_projection_surjective}.
    Hence $q\in\surj{\carrier[R]}{\carrier[\QuotientRing{R}{I}]}$ by
        \auto{additive_group_fields,quotient_ring_carrier}.
    Follows by
        \auto{ring_surjection,quotient_ring_projection_hom}.
\end{proof}

\begin{theorem}\label{quotient_ring_projection_kernel}
    Suppose $I$ is a two-sided ideal of $R$.
    Let $q=\groupQuotientProjection{\AdditiveGroup{R}}{I}$.
    Then $\ringKernel{q}{R}=I$.
\end{theorem}
\begin{proof}
    Let $G=\AdditiveGroup{R}$ and $Q=\quotientGroup{G}{I}$.
    $G$ is a commutative group by
        \obv{two_sided_ideal_ring,
            additive_group_is_exact_commutative_group}.
    $\subgroup{I}{G}$ by
        \auto{two_sided_ideal,left_ideal_additive_subgroup}.
    Hence $\normalSubgroup{I}{G}$ by
        \obv{commutative_group_subgroup_is_normal}.
    $\AdditiveGroup{\QuotientRing{R}{I}}=Q$ by
        \auto{quotient_ring_additive_group}.
    Follows by
        \auto{ring_kernel,group_quotient_projection_kernel}.
\end{proof}

\begin{theorem}\label{quotient_ring_zero_eq_one_iff}
    Suppose $I$ is a two-sided ideal of $R$.
    Then
        $\ringZero[\QuotientRing{R}{I}]
            =\ringOne[\QuotientRing{R}{I}]$ iff
        $\ringOne[R]\in I$.
\end{theorem}
\begin{proof}
    Let $S=\QuotientRing{R}{I}$ and
        $q=\groupQuotientProjection{\AdditiveGroup{R}}{I}$.
    $q$ is a ring homomorphism from $R$ to $S$ by
        \auto{quotient_ring_projection_hom}.
    $\ringKernel{q}{R}=I$ by
        \auto{quotient_ring_projection_kernel}.
    $q(\ringZero[R])=\ringZero[S]$ by \obv{ring_hom_zero}.
    $q(\ringOne[R])=\ringOne[S]$ by \obv{ring_hom_one}.
    Show that if $\ringZero[S]=\ringOne[S]$, then
        $\ringOne[R]\in I$.
    \begin{subproof}
        Assume $\ringZero[S]=\ringOne[S]$.
        $\ringOne[R]\in\carrier[R]$ by
            \obv{two_sided_ideal_ring,ring_one_type}.
        $q(\ringOne[R])=\ringZero[S]$.
        Hence $\ringOne[R]\in\ringKernel{q}{R}$ by
            \auto{ring_kernel_iff}.
    \end{subproof}
    Show that if $\ringOne[R]\in I$, then
        $\ringZero[S]=\ringOne[S]$.
    \begin{subproof}
        Assume $\ringOne[R]\in I$.
        Hence $\ringOne[R]\in\ringKernel{q}{R}$.
        Hence $q(\ringOne[R])=\ringZero[S]$ by
            \auto{ring_kernel_iff}.
    \end{subproof}
    Follows by \auto.
\end{proof}

\begin{theorem}\label{quotient_ring_trivial_iff}
    Suppose $I$ is a two-sided ideal of $R$.
    Let $S=\QuotientRing{R}{I}$.
    Then $\carrier[S]=\{\ringZero[S]\}$ iff $\ringOne[R]\in I$.
\end{theorem}
\begin{proof}
    $S$ is an exact ring by \auto{quotient_ring_is_exact}.
    Show that if $\carrier[S]=\{\ringZero[S]\}$, then
        $\ringOne[R]\in I$.
    \begin{subproof}
        Assume $\carrier[S]=\{\ringZero[S]\}$.
        $\ringOne[S]\in\carrier[S]$ by \obv{ring_one_type}.
        Hence $\ringOne[S]\in\{\ringZero[S]\}$.
        Hence $\ringOne[S]=\ringZero[S]$ by \auto{singleton_iff}.
        Follows by \auto{quotient_ring_zero_eq_one_iff}.
    \end{subproof}
    Show that if $\ringOne[R]\in I$, then
        $\carrier[S]=\{\ringZero[S]\}$.
    \begin{subproof}
        Assume $\ringOne[R]\in I$.
        Hence $\ringZero[S]=\ringOne[S]$ by
            \auto{quotient_ring_zero_eq_one_iff}.
        Show that $\carrier[S]\subseteq\{\ringZero[S]\}$.
        \begin{subproof}
            Show that for every $x\in\carrier[S]$ we have
                $x\in\{\ringZero[S]\}$.
            \begin{subproof}
                Fix $x\in\carrier[S]$.
                $x=\ringMul[S](\ringOne[S],x)$ by
                    \obv{ring_mul_one_left}.
                Hence $x=\ringMul[S](\ringZero[S],x)$.
                Hence $x=\ringZero[S]$ by \obv{ring_mul_zero_left}.
                Follows by \auto{singleton_iff}.
            \end{subproof}
            Follows by \obv{subseteq}.
        \end{subproof}
        $\{\ringZero[S]\}\subseteq\carrier[S]$ by
            \obv{ring_zero_type,singleton_subset_intro}.
        Follows by \obv{subseteq_antisymmetric}.
    \end{subproof}
    Follows by \auto.
\end{proof}

\begin{theorem}\label{quotient_ring_is_commutative}
    Let $R$ be a commutative ring.
    Suppose $I$ is a two-sided ideal of $R$.
    Then $\QuotientRing{R}{I}$ is a commutative ring.
\end{theorem}
\begin{proof}
    Let $G=\AdditiveGroup{R}$ and $E=\subgroupCosetRelation{G}{I}$.
    Let $S=\QuotientRing{R}{I}$.
    $S$ is an exact ring by \auto{quotient_ring_is_exact}.
    Show that for all $U,V\in\carrier[S]$ we have
        $\ringMul[S](U,V)=\ringMul[S](V,U)$.
    \begin{subproof}
        Fix $U,V\in\carrier[S]$.
        Take $a\in\carrier[G]$ such that
            $U=\equivalenceClass{E}{a}$ by
            \auto{quotient_ring_carrier,group_quotient_carrier,quotient}.
        Take $b\in\carrier[G]$ such that
            $V=\equivalenceClass{E}{b}$ by
            \auto{quotient_ring_carrier,group_quotient_carrier,quotient}.
        Hence $a,b\in\carrier[R]$ by \auto{additive_group_fields}.
        $\ringMul[S](U,V)=
            \equivalenceClass{E}{\ringMul[R](a,b)}$ and
            $\ringMul[S](V,U)=
            \equivalenceClass{E}{\ringMul[R](b,a)}$ by
            \auto{quotient_ring_mul_apply}.
        $\ringMul[R](a,b)=\ringMul[R](b,a)$ by
            \obv{commutative_ring_mul_commutative}.
    \end{subproof}
    Hence $S$ is a commutative ring by \obv{commutative_ring_intro}.
\end{proof}

\begin{corollary}\label{commutative_quotient_ring_is_commutative}
    Let $R$ be a commutative ring.
    Let $I$ be an ideal of $R$.
    Then $\QuotientRing{R}{I}$ is a commutative ring.
\end{corollary}
\begin{proof}
    $I$ is a two-sided ideal of $R$ by
        \auto{commutative_ideal_iff_two_sided_ideal}.
    Follows by \obv{quotient_ring_is_commutative}.
\end{proof}
